\documentclass[10pt, aspectratio=169]{beamer}

% --- Пакеты для русского языка и математики ---
\usepackage[T2A]{fontenc}
\usepackage[utf8]{inputenc}
\usepackage[russian]{babel}
\usepackage{amsmath, amsfonts, amssymb, amsthm}
\usepackage{graphicx}
\usepackage{tikz}

% --- МОДИФИКАЦИЯ ТЕМЫ ---
\usetheme{Madrid}
\useinnertheme{rounded} % Закругленные углы для блоков

% Настройка собственной цветовой схемы
\definecolor{DeepBlue}{RGB}{20, 40, 80}
\definecolor{ScienceGreen}{RGB}{0, 120, 80}

\setbeamercolor{palette primary}{bg=DeepBlue, fg=white}
\setbeamercolor{palette secondary}{bg=ScienceGreen, fg=white}
\setbeamercolor{palette tertiary}{bg=white, fg=DeepBlue}
\setbeamercolor{titlelike}{parent=palette primary}
\setbeamercolor{block title}{bg=DeepBlue, fg=white}
\setbeamercolor{block body}{bg=gray!10, fg=black}

% --- АВТОМАТИЧЕСКОЕ ОГЛАВЛЕНИЕ В НАЧАЛЕ СЕКЦИЙ ---
\AtBeginSection[]
{
  \begin{frame}{План выступления}
    \tableofcontents[currentsection]
  \end{frame}
}

% --- ТИТУЛЬНЫЙ ЛИСТ ---
\title[Теория графов]{Введение в спектральную теорию графов}
\subtitle{Методы и приложения в анализе данных}
\author[Павлов А.С.]{Павлов Александр Сергеевич}
\institute[ПЕТРГУ]{Петрозаводский государственный университет}
\date[2026]{Май 2026}

\begin{document}

% Слайд 1: Титульный
\begin{frame}
    \titlepage
\end{frame}

% Слайд 2: Содержание
\begin{frame}{Содержание презентации}
    \tableofcontents
\end{frame}

% --- СЕКЦИЯ 1: Базовые понятия ---
\section{Основные определения}

% Слайд 3: Определение
\begin{frame}{Математическая модель графа}
    \begin{definition}
        Неориентированным графом $G$ называется упорядоченная пара $(V, E)$, где:
        \begin{itemize}
            \item $V$ — непустое множество \textbf{вершин};
            \item $E$ — множество пар вершин, называемых \textbf{ребрами}.
        \end{itemize}
    \end{definition}
    \pause
    \begin{exampleblock}{Пример}
        Социальная сеть: пользователи ($V$) и связи между ними ($E$).
    \end{exampleblock}
\end{frame}

% Слайд 4: Матричное представление
\begin{frame}{Матричное представление графа}
    Для анализа структуры графа используется матрица смежности $A$:
    \begin{equation}
        A_{ij} = 
        \begin{cases} 
        1, & \text{если } (v_i, v_j) \in E \\
        0, & \text{в противном случае}
        \end{cases}
    \end{equation}
    \vspace{0.5cm}
    \begin{itemize}
        \item<2-> Матрица всегда симметрична для неориентированных графов.
        \item<3-> Сумма элементов $i$-й строки равна степени вершины $d(v_i)$.
    \end{itemize}
\end{frame}

% --- СЕКЦИЯ 2: Фундаментальные теоремы ---
\section{Спектральные свойства}

% Слайд 5: Теорема и нумерованная формула
\begin{frame}{Основная лемма о степенях}
    Рассмотрим связь между количеством ребер и степенями вершин.
    
    \begin{theorem}[Лемма о рукопожатиях]
        Для любого графа $G = (V, E)$ сумма степеней всех его вершин является четным числом и равна удвоенному количеству ребер.
    \end{theorem}
    
    \begin{equation}
        \sum_{v \in V} \text{deg}(v) = 2|E|
    \end{equation}
    \pause
    \textit{Следствие:} В любом графе число вершин нечетной степени четно.
\end{frame}

% Слайд 6: Оператор Лапласа
\begin{frame}{Матрица Лапласа}
    Важнейшим объектом спектральной теории является матрица $L = D - A$:
    
    \begin{equation}
        L_{ij} = 
        \begin{cases} 
        d(v_i), & \text{если } i = j \\
        -1, & \text{если } i \neq j \text{ и } (v_i, v_j) \in E \\
        0, & \text{в противном случае}
        \end{cases}
    \end{equation}
    
    \pause
    \begin{block}{Свойство}
        Спектр матрицы $L$ несет информацию о связности графа.
    \end{block}
\end{frame}

% --- СЕКЦИЯ 3: Интерактивные фрагменты ---
\section{Динамика и визуализация}

% Слайд 7: Поэтапный показ (Interactive 1)
\begin{frame}{Поэтапный алгоритм обхода}
    Процесс обработки графа в ширину (BFS):
    \begin{enumerate}
        \item<1-> Поместить начальную вершину в очередь.
        \item<2-> Пометить вершину как посещенную.
        \item<3-> Добавить всех соседей в конец очереди.
        \item<4-> Повторять, пока очередь не пуста.
    \end{enumerate}
    \uncover<5->{
        \begin{alertblock}{Важно!}
            Сложность алгоритма: $O(V + E)$.
        \end{alertblock}
    }
\end{frame}

% Слайд 8: Смена элементов (Interactive 2)
\begin{frame}{Сравнение типов графов}
    Нажмите кнопку «далее», чтобы увидеть отличия:
    \vspace{0.5cm}
    
    \centering
    \alt<1>{
        \textbf{Плотные графы} \\
        \begin{itemize}
            \item Число ребер $|E| \approx |V|^2$.
            \item Хранятся матрицами смежности.
        \end{itemize}
    }{
        \textbf{Разреженные графы} \\
        \begin{itemize}
            \item Число ребер $|E| \approx |V|$.
            \item Хранятся списками смежности.
        \end{itemize}
    }
\end{frame}

% Слайд 9: Динамические формулы (Interactive 3)
\begin{frame}{Эволюция спектра}
    Как меняется нормализованный Лапласиан?
    \vspace{1cm}
    \[
    \mathcal{L} = 
    \only<1>{D^{-1/2} L D^{-1/2}}
    \only<2>{I - D^{-1/2} A D^{-1/2}}
    \]
    \vspace{1cm}
    \begin{itemize}
        \item<1> Показывается стандартное определение.
        \item<2> Показывается форма через единичную матрицу $I$.
    \end{itemize}
\end{frame}

% Слайд 10: Заключение
\begin{frame}{Заключение}
    \begin{itemize}
        \item Теория графов — основа современного анализа сетей.
        \item Спектральные методы позволяют эффективно сжимать данные.
        \item \textbf{Спасибо за внимание!}
    \end{itemize}
    \vspace{1cm}
\end{frame}

\end{document}